\documentclass[12pt]{article}
\usepackage[margin=1in]{geometry}
\usepackage[T1]{fontenc}
\usepackage{bold-extra}
\usepackage{xcolor}
\definecolor{garnet}{HTML}{73000A}

\setlength{\parindent}{0pt}
\setlength{\parskip}{0.5em}

\begin{document}

{\color{garnet}\large\textbf{\textsc{SC Space Grant Personal Essay}}}\\[-0.8em]
{\color{garnet}\rule{\textwidth}{0.5pt}}\\[0.2em]

It should not have flown. A gyrocopter built from car transmissions, helicopter rotors, airplane controls, and snowmobile skis. The pilot at USC Fly Club had assembled it himself from parts never meant to share an airframe. But it flew.

The pilot did not need aerospace-grade components. He needed to understand the essential physics of flight and find parts that satisfied those requirements. That distinction has shaped how I approach every problem since.

Conversations at the Fly Club reinforced this idea. Pilots described how interconnected aircraft systems could be replaced or approximated by simpler mechanisms when the essential physics remained intact. Meanwhile, my coursework in aerospace engineering formalized the same intuition. Building powertrain simulations in MATLAB and Simscape taught me to choose appropriate abstractions. Not every detail needs modeling. The skill is knowing which ones do.

In Dr.\ Yi Wang's lab, I work on radar systems for autonomous maritime navigation. I have fine-tuned object detection models for deployment on embedded hardware, built pipelines that convert raw radar returns into usable imagery, and designed sensor mounts for research vessel integration. I present progress to Department of Defense sponsors monthly. Through this work, a persistent challenge emerged. Validating navigation algorithms requires real-world data, and real-world data requires expensive field campaigns. Testing an algorithm for conditions you cannot easily replicate means either waiting for the right opportunity or finding another way.

My proposed research addresses this gap directly. I am developing an open-source radar simulator that models terrain occlusion through geometric ray-casting rather than expensive electromagnetic computation. By validating the simulator against real offshore measurements, I will establish quantitative accuracy bounds. Researchers will know when the simplified model suffices and when higher fidelity is required. This makes algorithm validation accessible to teams without the budget for full physics-based simulation.

Before any autonomous system flies, its algorithms are tested thousands of times in simulation. NASA's Mars 2020 lander validated terrain-relative navigation this way, achieving 5-meter accuracy after extensive ground testing. The same principle applies to lunar rovers, autonomous aircraft, and future planetary explorers. Validated simulation tools are not optional infrastructure. They are how autonomous systems earn the trust required for flight.

The gyrocopter pilot did not wait for perfect parts. He built what he needed from what he had. I aim to do the same for radar simulation. Not a perfect electromagnetic solver, but a validated tool that works. Accessible, open-source, and ready for the researchers who need it.

\end{document}
